\documentclass[12pt]{article}
\nonstopmode                   \nobreak
\input ../input/colors
\input ../input/math
\topmargin 5mm                 \oddsidemargin 2mm
\textwidth 155mm               \textheight 245mm
\headheight 0mm                \headsep 0mm
\pagestyle{empty}
\parindent=0pt                 \parskip=15pt
\newcommand{\tw}[1]{\texttt{#1}}

\begin{document}\normalsize
\baselineskip=16pt plus 0.5pt

Astra version 6.1 enables running several (up to 20) parallel processes.
Operation is similar to a subroutine (plug-in) call in the previous Astra
versions. 
The only difference with calling a plug-in module is the ampersand
symbol, {\tt\&}, in the command line. 
For instance, the compiler instruction 

\hspace*{25pt}{\tt
NEUT~\&~:~0.1 :~2.2~:~3.3;
}\hfill(1)

will invoke the \tw{NEUT} code as a separate processes with the name
\tw{'neut'} inside the time interval $2.2\mbox{~s}\leq t\leq 3.3\mbox{~s}$. 
The process \tw{neut} will be activated and its results inquired every
$0.1\mbox{~s}$. 
Without the {\tt\&} symbol in the command line, the standard plug-in module
\tw{NEUT} will be called.
This new feature provides several advantages:
\vspace*{-\parskip}
\begin{enumerate}
\item It allows to call different subroutines or pieces of subroutines
  in parallel rather than in sequence as before.
  It opens an easy way to parallelizing the code and gaining from
  multi-processor computers. 
\item Maintenance and development of Astra and different routines become fully
  independent. 
  Change on one side does not require any awareness of all other sides as long
  as the interface does not change.
%  Unlike that, in the conventional approach, each modification of requires  
\item With growing complexity of the codes this approach allows to keep the
  responsibility of the code developers split and provides the only way of 
  maintaining huge code packages without increasing a number of people
  involved. 
\end{enumerate}\vspace*{-\parskip}

The IPC (inter-process communication) mechanism is used for synchronization
and communication with the Astra core. 
The synchronization is implemented by the Astra compiler and
does not require any additional programming on the user side.
Nevertheless, some module-specific features cannot be added automatically and
have to be provided by the user. 
Below we explain the code flow, processes interaction, data exchamge 
and then illustrate the aforesaid with an example.  

Assume that, in the previous Astra version, a subroutine \tw{example()} with
the source code \tw{sbr/example.c} has been used as a plug-in module. 
Transforming it to an independently running child process
requires a development of

\vspace*{-\parskip}
\begin{enumerate}
\item[1.1] multiple process synchronization (semaphore operations),
\item[1.2] two-sided shared memory read/write facilities (inter-processes
  communication),
\item[1.3] a stand-alone program \tw{main()} calling {~\tt example()} as a
  subroutine (function).
\end{enumerate}\vspace*{-\parskip}

As mentioned, the first step (1.1) is fully performed with the Astra compiler
and does not require any additional actions from the user.
To a large extent, the rest of the programming work is similar to including a
new plug-in C-module into the Astra code. 
In this case, one needs to
\vspace*{-\parskip}
\begin{enumerate}
\item[2.1] develop an interface for proper parameter exchange between
  \tw{example()} and the Astra Fortran core, 
\item[2.2] include a call of this interface into Astra.
\end{enumerate}\vspace*{-\parskip}

The step (2.1) consists of:
(i) identifying I/O parameter list and (ii) organizing exchange between the
Fortran and C pieces of the code.
Obviously, the step (i) is exactly the same in both cases (1.2) and (2.1). 
The step (ii) being formally also the same differs in implementation so far as
two processes communicate via shared memory.
The difference of two schemes can be represented as

\hskip5pt
\leftline{\tw{
Astra -> \color{blue}set\_example\_input -> example -> get\_example\_output
\color{black} -> Astra
}\hfill(2)~}

\hskip5pt
\leftline{\tw{
Astra -> \color{blue}set\_example\_input -> write\_Sh\_mem
}}\\
\centerline{\tw{\hfill
main -> \color{blue} read\_Sh\_mem -> example -> write\_Sh\_mem
}\hfill(3)}\\
\rightline{\tw{\color{blue}
read\_Sh\_mem -> get\_example\_output \color{black} -> Astra~~
}}

The two intermediate steps ~\tw{\color{blue}set\_example\_input}~ and
~\tw{\color{blue}get\_example\_output}~
are indispensable because the Astra core is written in Fortran while all
operations with the shared memory have to be written in C.
The program parts that are expected to be provided by the user are shown in
blue. 

The step (1.3) is shown as the central line in (3). 
In order to insure the proper time alignment it is recommended to 
embed the user developed piece into the pre-defined module
\tw{IPC/template.c}. 
The template module is supplied with the Astra code.
It contains a problem-independent environment that is responsible for
synchronizing with the Astra core and a specified place for including the
problem-dependent part (1.2). 
As shown in (3), the call of the C-function \tw{example()} is enclosed with two
series of commands supplying data to and receiving data from the function
\tw{example()}. 
Except for this operator sequence the file \tw{IPC/template.c} should remain
untouched.
The executable file must have the name \tw{IPC/example.c}. 
Instructions for compiling \tw{IPC/example.c} and preparing executable file 
\tw{.tsk/.../example} are given in the file \tw{IPC/Readme}. 

We describe now in more details some recommendations and constraints for
(1.2).  
First of all, a shared memory segment has to be designed that serves for
parameter exchange between the subprocess and the Astra core.
Note that, in the version 6.1, one and only one
shared memory segment has to be attached to each parallel process.
It is proposed to cast the exchange dataset in the form of the C-structure 
\\[5pt]
\hspace*{25pt}\parbox{10cm}{
\tw{struct A\_example\\
  \{\\
    struct A\_comio My;\\
    int  ...;\\
    double ...;\\
  \} *InOut;}}\\[5pt]
The structure is started with a service dataset
\\[2pt] 
\hspace*{25pt}\tw{struct A\_comio My;} 
%\hspace*{25pt}\parbox{10cm}{\tw
%struct A\_comio\\
%\{\\
%  char  Path[64];\\
%  key\_t Key;\\
%  pid\_t Pid;\\
%  int   OrdNr;\\
%  int   ShMid;\\
%  double CPUse;\\
%\};\\ }
\\[2pt] 
and then includes all Astra parameters required on the interface.
The user does not need to take care of the service dataset although it is the
user responsibility to the ensure a proper data exchange between the function 
\tw{example()} and the other data in the structure

\tw{struct A\_example *InOut};

In essence, the module (1.2) makes the same job as (2.1) excepting that the
assignment operators in (1.2) have to be replaced by communications with the  
shared memory.
Reading and writing shared memory includes altogether four modules.
In the child process, the reading and writing pieces can be placed
in the main file \tw{example.c} (see \tw{IPC/*.c}).
Unlike that, on the Astra side, each of the read and write interfaces has to
be split in two, C- and Fortran-, modules where the former one deals with the
shared memory while the latter with the Astra core.
The names of these four modules on the Astra side are predefined and derivable
from the name 
\tw{example}. 
In our case, they have to be named 

\vskip5pt\centerline{\hfill\tw{
toexample -> to\_example() \dots\ ot\_example() -> otexample
}\hfill(4)}

where the inner two are C-functions, the outer two Fortran subroutines.
The left hand side writes a specified set of Astra parameters to the shared
memory segment associated with the process \tw{example()} while the right hand
side reads the shared memory and transforms the output to Astra parameters.
Although the chain (4) could seem to be over-complicated it cannot be reduced
because it connects a C-function in the middle (dots) with the Fortran code
core on the edges.
Moreover, as long as the contents of all modules in (4) depend on the
particular subroutine \tw{example()} these functions cannot be created
automatically by the Astra compiler. 
Processing an instruction of the type (1) the Astra compiler adds call of two
subroutines \tw{toexample} and \tw{otexample}. 
Proper subroutines have to be provided by the user.

There are some further restrictions.
The Fortran subroutines \tw{toexample} and \tw{otexample} have no parameters.
They communicate with the Astra core via common blocks.
In opposite, the data transfer between \tw{toexample} and \tw{to\_example} 
(\tw{ot\_example} and \tw{otexample}) is organized through the C-function
headers. 
The two reciprocal functions \tw{to\_example()} and 
\tw{ot\_example()} can be found in the file \tw{ipc/io\_example.c}. 
The four modules (3) have to be placed into the Astra subdirectory \tw{ipc/}. 
Then they will be automatically linked to the Astra core by the Astra compiler.
We remind that the executable file for the child process \tw{.tsk/.../example}
has to be provided by the user.

In order to simplify the user work the Astra compiler creates, fills and
updates two segments of shared memory with main Astra variables and arrays
that are described in the include files \tw{for/A\_vars.h} and
\tw{for/A\_arrs.h}, respectively. 
Note that a code created by the compiler only writes these two shared memory
segments.
It never reads data from those. 
For instance, consider the model of warm neutrals \tw{NEUT}. 
Although the neutral density \tw{NN} and temperature \tw{TN} are included
into the shared memory segment associated with \tw{for/A\_arrs.h} it is the
user responsibility to update these quantities after calling \tw{NEUT}.

In summary, we give below a list of new files and functions that have to be
created instead of the subroutine file \tw{sbr/example.c}

\vspace*{-\parskip}
\begin{center}
\begin{tabular}{|l|l|l|}
\hline
File & Function/Subroutine & Action \\
\hline
\tw{ipc/toexample.f}   & \tw{toexample}
                 & prepare Astra data for \tw{to\_example}\\
\tw{ipc/to\_example.c} & \tw{to\_example(\dots)} 
                 & write shared memory before \tw{example} \\
\hline
\tw{IPC/example.c} & \tw{example()} 
                         & process synchronization      \\
                 & 
                 & read shared memory before \tw{example}   \\
                 & %\tw{a\_example();} 
                 & call \tw{example}       \\
                 & %\tw{write\_shmem} 
                 & write shared memory after \tw{example}    \\
\hline
\tw{ipc/ot\_example.c} & \tw{ot\_example(\dots)} 
                 & read shared memory after \tw{example}   \\
\tw{ipc/otexample.f}   & \tw{otexample}
                 & transfer results of \tw{ot\_example} to Astra\\
\hline
\end{tabular}
\end{center}
Note that the names \tw{toexample} and \tw{otexample} are mandatory while all
other names in the first and second columns are optional.
For instance, one can put the two Fortran subroutines, \tw{toexample} and
\tw{otexample}, in one file \tw{ottoexample.f},  etc.

The situation becomes a little bit more complicated if one needs to convert a
Fortran subroutine \tw{sbr/example.f} to an independent process.
Then it is necessary to add one more interface between the child \tw{main}
C-function and the Fortran subroutine \tw{example.f}.
The procedure is standard, we only note that one can keep the
original source \tw{sbr/example.f} unchanged and call it in parallel with the
newly created process so long as a special care is taken that the outputs of
the two do not destroy one another.

\vfill\eject
\centerline{\bf Code flow diagram}
\begin{picture}(400,500)
%\put(110,380){\begin{picture}(50,50)%
\put(290,-130){\begin{picture}(50,50)%
                \put(-10,30){\fbox{\fbox{\parbox{200pt}{\begin{center}\bf
                        \color{ippblue}
                        Astra initiates semaphores, 
                        launches all secondary processes, 
                        inquires them, then locks
                        \color{black}\end{center}}}}}
             \end{picture}}
                \put(155,470){\fbox{\parbox{115pt}{\begin{center}
                        \color{ippblue}
                        \bf Create semaphores%\vphantom{y}
                        \color{black}\end{center}}}}
                \put(155,410){\fbox{\parbox{115pt}{\begin{center}
                        \color{ippblue}
                        \bf Launch secondaries
                            Lock primary%\vphantom{y}
                        \color{black}\end{center}}}}
                \put(105,379){\line(1,0){195}}
                \put(300,370){\line(0,1){90}}
                \put(30,375){\parbox{105pt}{\bf\color{red}
                        Checkpoint
                        \color{black}}}
    \put(155,80){\begin{picture}(50,50)%
                \put(160,400){\fbox{\parbox{105pt}{\begin{center}
                        \color{ippblue}
                        \bf Check existance\vphantom{y}
                        \color{black}\end{center}}}}
                \put(160,340){\fbox{\parbox{105pt}{\begin{center}
                        \color{ippblue}
                        \bf Collect info\vphantom{y}
                        \color{black}\end{center}}}}
                \put(160,280){\fbox{\parbox{105pt}{\begin{center}
                        \color{ippblue}
                        \bf Create shmem\vphantom{y}
                        \color{black}\end{center}}}}
             \end{picture}}

                \put(160,350){\fbox{\parbox{105pt}{\begin{center}
                        \color{ippblue}
                        \bf Wait for secs%\vphantom{y}
                        \color{black}\end{center}}}}
                \put(160,290){\fbox{\parbox{105pt}{\begin{center}
                        \color{ippblue}
                        \bf Get shmem requisites%\vphantom{y}
                        \color{black}\end{center}}}}

\put(0,-60){\begin{picture}(50,50)%\red%
        \put(210,-58){\line(0,-1){7}}
        \put(210,-65){\vector(-1,0){105}}
        \put(105,-65){\line(-1,0){105}}
        \put(0,-65){\vector(0,1){190}}
        \put(0,125){\line(0,1){190}}
        \put(0,315){\vector(1,0){105}}
        \put(105,315){\line(1,0){105}}
        \put(210,315){\line(0,-1){10}}
             \end{picture}}
                \put(160,220){\fbox{\parbox{105pt}{\begin{center}
                        \color{ippblue}
                        \bf Write shmem\vphantom{y}
                        \color{black}\end{center}}}}
                \put(155,150){\fbox{\parbox{115pt}{\begin{center}
                        \color{ippblue}
                        \bf Unlock secondaries Lock primary\vphantom{y}
                        \color{black}\end{center}}}}
                \put(20,115){\line(1,0){280}}
                \put(20,52){\line(1,0){280}}
                \put(300,115){\line(0,1){80}}
                \put(300,52){\line(0,-1){90}}
                \put(300,195){\line(1,0){120}}
                \put(300,-38){\line(1,0){120}}
                \put(20,80){\fbox{\parbox{105pt}{\begin{center}
                        \color{ippblue}
                        \bf Run secondary\vphantom{y}
                        \color{black}\end{center}}}}
                \put(160,20){\fbox{\parbox{105pt}{\begin{center}
                        \color{ippblue}
                        \bf Wait for secs\vphantom{y}
                        \color{black}\end{center}}}}
                \put(160,-40){\fbox{\parbox{105pt}{\begin{center}
                        \color{ippblue}
                        \bf Read shmem\vphantom{y}
                        \color{black}\end{center}}}}
                \put(160,-100){\fbox{\parbox{105pt}{\begin{center}
                        \color{ippblue}
                        \bf Time step\vphantom{y}
                        \color{black}\end{center}}}}
    \put(155,80){\begin{picture}(50,50)%
                \put(5,0){\fbox{\parbox{105pt}{\begin{center}
                        \color{ippblue}
                        \bf Run secondary\vphantom{y}
                        \color{black}\end{center}}}}
                \put(-130,-37){\parbox{105pt}{\bf\color{red}
                        Checkpoint
                        \color{black}}}
                \put(-120,40){\parbox{105pt}{\bf\color{red}
%                        1st checkpoint
                        \color{black}}}
                \put(-120,-40){\parbox{105pt}{\bf\color{red}
%                        2nd checkpoint
                        \color{black}}}
                \put(170,120){\parbox{105pt}{\bf\color{red}
%                        1st checkpoint
                        \color{black}}}
                \put(170,-130){\parbox{105pt}{\bf\color{red}
%                        2nd checkpoint
                        \color{black}}}
                \put(160,80){\fbox{\parbox{105pt}{\begin{center}
                        \color{ippblue}
                        \bf Read shmem\vphantom{y}
                        \color{black}\end{center}}}}
                \put(160,30){\fbox{\parbox{105pt}{\begin{center}
                        \color{ippblue}
                        \bf Compute\vphantom{y}
                        \color{black}\end{center}}}}
                \put(160,-20){\fbox{\parbox{105pt}{\begin{center}
                        \color{ippblue}
                        \bf Write shmem\vphantom{y}
                        \color{black}\end{center}}}}
                \put(160,-80){\fbox{\parbox{105pt}{\begin{center}
                        \color{ippblue}
                        \bf Inform primary Lock itself\vphantom{y}
                        \color{black}\end{center}}}}
             \end{picture}}
\end{picture}
\vfill\eject

\centerline{\bf\large Code talking mechanism}

\begin{enumerate}
\item At start, a set of~ {\blue\tt Nsbp+1} ~semaphores is created.\\ 
  All semaphores in the set are set to zero.
\item {\blue\tt Nsbp+2} ~shared memory segments are created.
\item The IDs of the primary {\blue\tt main} and of all secondary processes 
  as well as semaphore and shared memory IDs are stored in the file 
  {\blue\tt AWD/tmp/astra.ipc}.
\end{enumerate}


\leftline{\bf Notations}\\
The box
\begin{picture}(70,0)%
   \put(10,0){\fbox{\fbox{\parbox{40pt}{\bf
       \color{ippblue}~~\# i\color{black}}}}}
\end{picture}
means that\vspace*{-1em}
\begin{enumerate}
\item On entry, {\blue\tt sem\_val(i)=0}. This causes $i$-th subprocess
  {\blue\tt SBP(i)} to wait until another process ({\blue\tt main} in case
  {\blue\tt i$\neq$0} or an SBP otherwise) increments the $i$-th semaphore
  by setting its value to {\blue\tt sem\_val(i)=1}. 
\item Then {\blue\tt SBP(i)} proceeds simultaneously resetting the semaphore
  value back to zero {\blue\tt sem\_val(i)=0}. This restores the initiial
  status and will lock the process {\blue\tt SBP(i)} on the next entry. 
\end{enumerate}

\vskip1em
\leftline{\sc \underline{Start}}
{\bf Initial status:}\\ All semaphores are set to zero:
{\blue\tt sem\_val[0:Nsbp]=0}. \\[15pt]
{\blue\tt inikids:\hspace*{16em}SBP(i)}\\
\hspace*{5em}
\begin{picture}(300,155)%
   \put(30,160){\dots}
   \put(20,140){\blue\tt i = 1;}
   \put(0,120){\rm Launch~~ {\tt SBP(i)}}
   \put(15,90){\fbox{\fbox{\parbox{40pt}{\bf
       \color{ippblue}~~\# 0\color{black}}}}}
   \put(-20,45){No}
   \put( 10,40){\line(-1,0){50}}
   \put(-40,40){\line( 0,1){85}}
   \put(-40,125){\vector( 1,0){30}}
   \put(35, 30){\vector(0,-1){35}}
   \put(40,10){Yes}
   \put(25,60){\blue\tt i++;}
   \put(20,40){\em Ready?}
                        \put(170,120){\rm Collect IPC info}
                        \put(150,100){\rm Write~1~line~to~\tt A\_ipc\_file}
                        \put(170,80){\blue\tt sem\_val(0)++;}
                        \put(180,55){\fbox{\fbox{\parbox{40pt}{\bf
                                \color{ippblue}~~\# i\color{black}}}}}
                        \color{ippblue}
                \put(100,125){\vector(1,0){60}}
                \put(160,80){\vector(-1,0){80}}
                \put(100,85){Release}
                \put(105,70){main}
                        \color{black}
\end{picture}

%\vskip10pt
{\bf Resulting status:} \\[5pt]
\quad All semaphores are set to zero $~\longrightarrow~$ All SBPs wait.\\
{\blue\tt main} ~runs until the next checkpoint (see main Loop).

\vfill\eject
\leftline{\sc \underline{Main Loop}}
{\bf Status.} All semaphores are set to zero:
{\blue\tt sem\_val[0:Nsbp]=0}. 
All SBPs are sleeping.

{\blue\tt main:\hspace*{26em}SBP(i)}\\
\hspace*{6em}
\begin{picture}(300,300)%
   \put(30,305){\dots}
   \put(-15,285){\blue\tt i = 1;~~buf0.sem\_op = 0;}
   \put( 0,265){\rm Check if~ {\tt SBP(i)} is due?}
   \put(20,255){\vector(0,-1){110}}
   \put(25,240){No}
   \put(85,255){\vector(0,-1){25}}
   \put(90,240){Yes}
   \put(50,215){\blue\tt buf0.sem\_op$\,$-$\,$-;}
   \put(50,200){\blue\tt sem\_val(i)++;}
%   \put(60,190){\fbox{\fbox{\parbox{40pt}{\bf
%       \color{ippblue}~~\# i\color{black}}}}}
   \put(85,170){\vector(0,-1){25}}
   \put(-20,130){\rm Check if~ {\tt SBP(i)} list is exausted?}
   \put( 85,120){\vector(0,-1){50}}
   \put(90,90){Yes}
   \put(20,120){\vector(0,-1){25}}
   \put(25,105){No}
   \put(10, 80){\blue\tt i++;}
   \put( 5, 80){\line(-1,0){45}}
   \put(-40,80){\line( 0,1){190}}
   \put(-40,270){\vector( 1,0){30}}
   \put(-40,50){Check$\red^*$:~{\blue\tt sem\_val(0)$\geq$ abs(buf0.sem\_op)}}
   \put( 85,40){\vector(0,-1){55}}
   \put(90,10){Yes}
   \put(20,40){\vector(0,-1){25}}
   \put(25,25){No}
   \put(5, 0){Wait\red$^{**}$}
   \put(0, 0){\line(-1,0){60}}
   \put(-60,0){\line(0,1){55}}
   \put(-60,55){\vector(1,0){15}}
   \put(0,-30){Make the next time step}
   \put(5,-48){Repeat the loop}
   \put(170,0){\red Time synchronization checkpoint}
   \put(-50,-5){\red \rule{400pt}{1pt}}
   \put(0,-45){\line(-1,0){70}}
   \put(-70,-45){\line(0,1){340}}
   \put(-70,295){\vector(1,0){105}}
        \color{ippblue}
        \put(135,202){\vector(1,0){120}}
        \put(170,207){Release}
        \put(172,189){SBP(i)}
        \color{black}
                \put(260,200){\fbox{\fbox{\parbox{40pt}{\bf
                        \color{ippblue}~~\# i\color{black}}}}}
                \put(250,175){Makes its job}
                \put(250,160){\blue\tt sem\_val(0)++;}
                \put(290,155){\line(1,0){50}}
                \put(340,155){\line(0,1){70}}
                \put(340,225){\vector(-1,0){50}}
\end{picture}

\vskip40pt
{\bf Note}
\begin{enumerate}
\item \vskip-20pt
      Subroutines and subprocesses are launched sequentially in the same order
      as they are listed in a model file unless the tag {\tt "<"} is used.
\item Moreover the subroutines are executed in sequence hence the next one is
  not started before the previous one finishes its calculation.
\item Checkpoint is placed after all subroutines and subprocesses before the
  first time evolution element ({\tt ADCMP}).
\end{enumerate}
\vskip-20pt
{\red \rule[2pt]{2cm}{1pt}}
 ~

\vfill
\rule{80pt}{0.5pt}\\
$\red^*$~Here {\blue\tt sem\_val(0)} is a number of SBPs that have finished
their work and are waiting, {\blue\tt $-$(buf0.sem\_op)} is a number of SBPs
that have been called at the current time step.\\
$\red^{**}$~Here {\blue\tt main} can do also some operations (e.g. polling
events, plotting, etc) that are not directly involved in the time evolution.\\

\end{document}
   \put(20,40){\em Ready?}
                        \put(170,120){\rm Collect IPC info}
                        \put(150,100){\rm Write~1~line~to~\tt A\_ipc\_file}
                        \put(170,80){\blue\tt sem\_val(0)++;}
                        \put(180,55){\fbox{\fbox{\parbox{40pt}{\bf
                                \color{ippblue}~~\# i\color{black}}}}}
                        \color{ippblue}
                \put(160,80){\vector(-1,0){80}}
                \put(100,85){Release}
                \put(105,70){main}
                        \color{black}

